\section{Heights}

We now introduce the height of a commutative formal group law.  In this
section, we fix a field $K$ of characteristic $p$, and we consider a
commutative formal group law \lstinline|F| over $K$.

The first step is to define the multiplication-by-\lstinline|n| series of
the formal group law.  Since we have already constructed an
\lstinline|AddCommGroup| instance on \lstinline|F.Point σ| under the
assumption \lstinline|F.IsComm|, we can directly use the scalar
multiplication notation \lstinline|n • x| for repeated addition of a
point.

\begin{lstlisting}
def series (n : ℕ) : FormalGroupHom F F where
  toPowerSeries := (n • ⟨PowerSeries.X, PowerSeries.HasSubst.X _⟩ : F.Point Unit).val
  zero_constantCoeff := ...
  hom := ...
\end{lstlisting}

Here \lstinline|⟨PowerSeries.X, ...⟩| is the point of
\lstinline|F.Point Unit| given by the coordinate power series
\lstinline|PowerSeries.X|, and \lstinline|n • ⟨PowerSeries.X, ...⟩| is its
\lstinline|n|-fold sum under the addition induced by \lstinline|F|; taking
\lstinline|.val| extracts the underlying power series.  Thus
\lstinline|F.series n| is the formal group homomorphism for multiplication
by \lstinline|n|, usually denoted $[n]_F(X)$.  The field
\lstinline|zero_constantCoeff| shows it has constant coefficient zero, and
\lstinline|hom| proves compatibility with the formal group law structure;
we omit these proofs.

The definition of height focuses on the multiplication-by-\lstinline|p|
series \lstinline|F.series p|.  The following theorem shows that, if it is
nonzero, then it has a nonzero coefficient in some degree of the form
\lstinline|p ^ n|.

\begin{lstlisting}
theorem exist_coeff_pow_ne_zero_iff_ne_zero :
    (∃ n, (F.series p).toPowerSeries.coeff (p ^ n) ≠ 0) ↔
      (F.series p).toPowerSeries ≠ 0 := ...
\end{lstlisting}

Using the alternative definition in Remark~\ref{rem: alter_def_height}, when \lstinline|F.series p| is nonzero, it
yields a smallest \lstinline|n| with a nonzero coefficient in degree
\lstinline|p ^ n|. The height of \lstinline|F| is then defined as follows:
\begin{lstlisting}
def height : ℕ∞ :=
  letI := Classical.decEq K
  letI := Classical.decEq (PowerSeries K)
  if h : (F.series p).toPowerSeries = 0 then ⊤ else
    Nat.find ((F.exist_coeff_pow_ne_zero_iff_ne_zero).mpr h)
\end{lstlisting}

The value of \lstinline|height| lies in \lstinline|ℕ∞|, the natural
numbers together with a top element \lstinline|⊤|.  If
\lstinline|F.series p| is zero, the height is \lstinline|⊤|; otherwise
\lstinline|Nat.find| chooses the smallest \lstinline|n| with a nonzero
coefficient in degree \lstinline|p ^ n|.

In ordinary notation, the height records the first degree of the form
\lstinline|p ^ n| in which \lstinline|[p]_F(X)| has a nonzero coefficient:
\lstinline|F| has height \lstinline|n| if
\[
  [p]_F(X) = a X^{p^n} + \text{higher order terms}
\]
with \lstinline|a| nonzero, and infinite height if
\lstinline|[p]_F(X) = 0|.